\documentclass[12pt]{article}
\usepackage[a4paper,margin=1in]{geometry}
\usepackage[T1]{fontenc}
\usepackage{lmodern,microtype}
\usepackage{amsmath,amssymb,amsthm,mathtools}
\usepackage{booktabs,array,enumitem,xcolor}
\usepackage[numbers,sort&compress]{natbib}
\usepackage[colorlinks=true,linkcolor=blue!55!black,citecolor=blue!55!black,urlcolor=blue!55!black]{hyperref}
\linespread{1.07}

\newtheorem{theorem}{Theorem}[section]
\newtheorem{proposition}[theorem]{Proposition}
\newtheorem{corollary}[theorem]{Corollary}
\theoremstyle{remark}
\newtheorem{remark}[theorem]{Remark}
\newcommand{\norm}[1]{\left\lVert#1\right\rVert}

\title{Regularized Cross-Gram Tradeoffs\\
\large Singular-Value Clipping Cannot Beat a Residual Larger Than the Cross Angle}
\author{Liang Wang}
\date{July 2026}

\begin{document}
\maketitle

\begin{abstract}
The RH-185 bi-Krylov candidate has small directed residuals only in windows
where the right and left temporal spaces are nearly transverse.  RH-186
shows that the elementary oblique conditioning gate fails.  A natural repair
is to clip the small singular values of the cross Gram and accept a controlled
loss of biorthogonality.  This paper derives the exact tradeoff.

Let $\gamma=\sigma_{\min}(Q_L^*Q_R)$, let $\epsilon$ be the larger directed
residual, and clip at $\tau>0$.  The effective singular value is
$d=\max(\gamma,\tau)$.  The regularized duality ratio is $\gamma/d$, the
duality defect is $1-\gamma/d$, and the regularized coordinate norm is
$1/d$.  Their sum obeys the exact identity
\begin{equation*}
 \left(1-\frac{\gamma}{d}\right)+\frac{\epsilon}{d}
 =1+\frac{\epsilon-\gamma}{d}
 \quad (\tau\ge\gamma).
\end{equation*}
Consequently a strict combined contraction exists in this clipping family if
and only if $\epsilon<\gamma$.  If the residual is larger than the cross
angle, clipping can trade conditioning for duality but cannot push the total
budget below one.

We sweep 13 clipping thresholds over all 126 RH-185 windows, producing 1,638
records.  The smallest residual-to-cross-angle ratio is $10.253$; no window
has a strict regularized contraction, and the best finite grid gate is
$1.0090$.  This is a sharp branch-specific negative result, not a theorem
excluding non-spectral regularization or contour cancellation.
\end{abstract}

\section{Motivation}

For balanced frames $V,W$ with cross-angle cosine $\gamma$, the natural
condition number is $\gamma^{-1}$ \citep{WangRH184,WangRH186}.  The twelve
RH-185 windows passing the raw residual gate have residual levels much larger
than their smallest cross singular values.  One might hope to replace the
exact dual frame by a bounded approximate dual.

Singular-value clipping is the cleanest target-independent test.  It acts on
the cross Gram itself, not on fitted eigenvectors.  If even this family cannot
meet a combined duality/conditioning gate, the next regularization must use
additional structure rather than merely truncating small cross angles.

\section{Clipped frames}

Write
\begin{equation}\label{eq:cross-svd}
 H=Q_L^*Q_R=U\operatorname{diag}(\sigma_1,\ldots,\sigma_r)V^*,
 \qquad \sigma_r=\gamma>0.
\end{equation}
For a clipping level $\tau>0$, set
\begin{equation}\label{eq:effective}
 d_j=\max(\sigma_j,\tau).
\end{equation}
Use the same singular vectors but replace $\sigma_j^{-1/2}$ by
$d_j^{-1/2}$:
\begin{equation}\label{eq:clipped-frames}
 V_\tau=Q_RV\operatorname{diag}(d_j^{-1/2}),
 \qquad
 W_\tau=Q_LU\operatorname{diag}(d_j^{-1/2}).
\end{equation}
Then
\begin{equation}\label{eq:approx-dual}
 W_\tau^*V_\tau
 =\operatorname{diag}\left(\frac{\sigma_j}{d_j}\right).
\end{equation}

The smallest singular direction controls the worst duality defect and the
largest clipped inverse controls the coordinate norm.  The following theorem
is the exact scalar statement needed for the audit.

\section{The clipping tradeoff theorem}

\begin{theorem}[Exact regularized budget]\label{thm:tradeoff}
Let $\gamma=\sigma_r$ and let $\epsilon\ge0$ be a residual level.  For the
clipped pair \eqref{eq:clipped-frames}, define
\begin{align}
 \delta_\tau&=1-\frac{\gamma}{\max(\gamma,\tau)},
 \label{eq:duality-defect}\\
 \chi_\tau&=\frac1{\max(\gamma,\tau)},
 \label{eq:regularized-condition}\\
 B_\tau&=\delta_\tau+\epsilon\chi_\tau.
 \label{eq:combined}
\end{align}
If $\tau\ge\gamma$, then
\begin{equation}\label{eq:identity}
 B_\tau=1+\frac{\epsilon-\gamma}{\tau}.
\end{equation}
There exists $\tau>0$ with $B_\tau<1$ if and only if $\epsilon<\gamma$.
\end{theorem}

\begin{proof}
For $\tau\le\gamma$, the pair is unaltered in the worst direction and
$B_\tau=\epsilon/\gamma$.  For $\tau\ge\gamma$, substitute
$\delta_\tau=1-\gamma/\tau$ and $\chi_\tau=1/\tau$ to obtain
\eqref{eq:identity}.  If $\epsilon<\gamma$, the unaltered choice gives
$B_\tau=\epsilon/\gamma<1$.  If $\epsilon=\gamma$, the value is one for
$\tau\le\gamma$ and approaches one from above for $\tau>\gamma$.  If
$\epsilon>\gamma$, the unaltered value exceeds one and
\eqref{eq:identity} is strictly larger than one.  Thus the stated
equivalence holds.
\end{proof}

\begin{corollary}[Cross-angle criterion]\label{cor:criterion}
Within singular-value clipping, regularization closes the combined
duality/conditioning gate exactly when
\begin{equation}\label{eq:criterion}
 \epsilon<\sigma_{\min}(Q_L^*Q_R).
\end{equation}
\end{corollary}

The result is a Pareto identity, not a numerical heuristic.  It says that
clipping cannot create transversality that the residual has already lost.

The same conclusion persists when the residual information is resolved in
the singular directions rather than replaced by one maximum.

\begin{proposition}[Directionwise clipping criterion]\label{prop:directionwise}
Let $\epsilon_j\ge0$ be a residual bound in the $j$th cross-Gram singular
direction and set $d_j=\max(\sigma_j,\tau)$.  The directionwise additive
budget
\begin{equation}\label{eq:directionwise-budget}
 B_j(\tau)=1-\frac{\sigma_j}{d_j}+\frac{\epsilon_j}{d_j}
\end{equation}
is strictly below one for every $j$ at some clipping level $\tau>0$ if and
only if
\begin{equation}\label{eq:directionwise-criterion}
 \epsilon_j<\sigma_j\qquad\text{for every }j.
\end{equation}
\end{proposition}

\begin{proof}
If \eqref{eq:directionwise-criterion} holds, choose
$\tau\le\sigma_r$.  Then no direction is clipped and
$B_j=\epsilon_j/\sigma_j<1$.  Conversely, if
$\epsilon_k\ge\sigma_k$ for some $k$, then for
$\tau\le\sigma_k$ one has $B_k=\epsilon_k/\sigma_k\ge1$, while for
$\tau>\sigma_k$ one has
$B_k=1+(\epsilon_k-\sigma_k)/\tau\ge1$.
\end{proof}

Thus the worst-direction theorem is not an artifact of collapsing all
singular vectors into one scalar.  Direction-dependent clipping can help
only when every affected residual direction is already smaller than its
available cross angle.

\section{Physical sweep}

For each RH-185 window we set
\begin{equation}\label{eq:epsilon}
 \epsilon=\max(\epsilon_R,\epsilon_L)
\end{equation}
and evaluate 13 thresholds
$\tau/\gamma\in\{0.25,0.5,1,2,\ldots,1024\}$.  This gives 1,638
threshold/window records.  The key values are:
\begin{center}
\begin{tabular}{lr}
\toprule
quantity & value\\
\midrule
windows & 126\\
thresholds per window & 13\\
minimum $\epsilon/\gamma$ & $10.2530$\\
median $\epsilon/\gamma$ & $902.203$\\
maximum $\epsilon/\gamma$ & $1.5855\times10^8$\\
minimum finite-grid $B_\tau$ & $1.00904$\\
strict regularized gates & 0\\
\bottomrule
\end{tabular}
\end{center}

The twelve RH-185 raw gate-passing windows are not exceptions: their
residual-to-cross-angle ratios are all larger than one.  In particular the
best right-channel $L=4$ window has residual level about $0.0419$ and
$\gamma\approx0.00104$, giving a ratio near $40$.  The best overall ratio,
$10.253$, still misses the exact criterion by an order of magnitude.

\section{What the negative result means}

The theorem rules out one very specific repair:
\begin{quote}\label{eq:branch-no-go}
Clip the cross-Gram singular values and pay for the lost duality in a single
additive budget.
\end{quote}
It does not rule out:
\begin{itemize}[leftmargin=2em]
\item a contour resolvent that exploits nonnormal cancellation;
\item a directed Schur product using separate left and right couplings;
\item a different physical dual seed or delayed start;
\item a regularizer based on the dynamics rather than cross-Gram singular
 values alone.
\end{itemize}
The next paper therefore studies the directed couplings before inserting a
complement resolvent.

\section{A useful diagnostic identity}

For $\tau\ge\gamma$, equation \eqref{eq:identity} can be rearranged as
\begin{equation}\label{eq:diagnostic}
 B_\tau-1=\frac{\epsilon-\gamma}{\tau}.
\end{equation}
The sign of every regularized gate is already decided by the unregularized
residual-to-angle comparison.  A sweep is still useful for quantifying the
distance to one and for detecting coding errors, but it cannot change that
sign.

\section{Pareto geometry and limiting regimes}

For $\tau\ge\gamma$, the two contributions move in opposite directions:
\begin{equation}\label{eq:pareto-components}
 \delta_\tau=1-\frac\gamma\tau\nearrow1,
 \qquad
 \epsilon\chi_\tau=\frac\epsilon\tau\searrow0.
\end{equation}
If $\epsilon>\gamma$, their sum decreases monotonically toward one from
above.  This explains why a sufficiently large clipping threshold can make
a finite sweep appear numerically close to success without ever producing a
strict certificate.  If $\epsilon<\gamma$, the unregularized value is
already below one and additional clipping increases the sum toward one.  At
$\epsilon=\gamma$, the sum is exactly one for every threshold.

Accordingly there is no hidden optimum to locate numerically.  The only
substantive datum is the sign of $\epsilon-\gamma$; the threshold controls
how the same failure or success is displayed.  In the physical sweep the
best finite-grid value $1.00904$ is therefore consistent with, rather than
contrary to, the exact no-contraction theorem.

\section{Reproducibility and branch scope}

The archive stores the original cross angle and residual level together
with all thirteen thresholds, effective singular values, duality defects,
regularized norm products, and combined gates.  Recomputing the sign from
the original pair $(\gamma,\epsilon)$ independently of the sweep is an
internal consistency check supplied by Theorem~\ref{thm:tradeoff}.

The negative result is intentionally modular.  It can be reused whenever a
future packet proposal suggests the same clipping repair: first compare its
residual with its cross angle.  If the comparison fails, no threshold search
in this additive model is needed.  If it passes, the unregularized balanced
pair already supplies the smallest gate in the family.

\section{Boundary}

This paper proves the exact clipping tradeoff and verifies it over 1,638
finite physical records.  It does not prove that all possible regularizers
fail, nor does it prove an all-level cross-angle obstruction.  The directional
Schur product, exact Feshbach factorization, and contour-specific complement
resolvent remain open.  R, Gate A, Hilbert--Polya, and RH remain untouched.

\bibliographystyle{plainnat}
\bibliography{references}
\end{document}
